\section{Notebook-Centric Methodology and Results}
This section documents the major notebook workflows in the repository for reproducible inclusion in a Master's thesis on speech emotion recognition.
\subsection{analysis/}
\textbf{What was done.} Exploratory notebooks (\texttt{iemocap\_eda.ipynb}, \texttt{mfcc\_output\_analysis.ipynb}) were used to inspect dataset composition and initial feature behavior.
\textbf{How it was done.} The notebooks are interactive and diagnostic; they focus on plots and sanity checks rather than a fixed train/test protocol.
\textbf{Results.} No standardized machine-readable result artifacts are currently exported from this folder.
\subsection{feature\_extraction/feature\_families/}
\textbf{What was done.} Multiple feature families were extracted from audio/text, including MFCC, prosody, tonality, rhythm, SSL embeddings, BERT embeddings, and TF-IDF.
\textbf{How it was done.} Family-specific notebooks/scripts compute per-utterance descriptors and export CSV files under \texttt{extracted\_features/}.
\textbf{Results.} Table~\ref{tab:feature_inventory} summarizes sample size and dimensionality per family.
\begin{table}[h]
\centering
\caption{Extracted feature family inventory}
\label{tab:feature_inventory}
\begin{tabular}{lrr}
\hline
Feature family & Samples & Feature count \\
\hline
BERT text embeddings & 7532 & 768 \\
MFCC (normalized) & 9887 & 1150 \\
MFCC (raw) & 9887 & 1150 \\
Prosody energy & 7532 & 8 \\
Prosody pitch & 7532 & 14 \\
Representations & 7532 & 136 \\
Rhythm/pauses & 7532 & 10 \\
SSL HuBERT & 9887 & 1539 \\
SSL wav2vec2 & 9887 & 1539 \\
TF-IDF text & 4779 & 10480 \\
Tonality & 7532 & 91 \\
\hline
\end{tabular}
\end{table}
\subsection{feature\_extraction/pipelines/}
\textbf{What was done.} Pipeline notebooks aggregate per-family features into combined datasets for downstream modeling.
\textbf{How it was done.} \texttt{all\_feature\_families\_pipeline.ipynb} and \texttt{anova\_selected\_feature\_pipeline.ipynb} orchestrate merges and dataset preparation.
\textbf{Results.} Current exported combined tables are:
\begin{table}[h]
\centering
\caption{Current pipeline outputs in extracted\_features/combined}
\begin{tabular}{lrr}
\hline
Output table & Rows & Columns \\
\hline
all\_families\_features & 7532 & 14034 \\
\hline
\end{tabular}
\end{table}
\subsection{feature\_selection/anova/select\_k\_best\_anova/}
\textbf{What was done.} ANOVA filter-based feature selection was evaluated per feature family.
\textbf{How it was done.} Each notebook runs baseline and post-selection training with class-imbalance-aware models (Logistic Regression, Linear SVC, Random Forest), then computes metric deltas.
\textbf{Results.} Table~\ref{tab:anova_results} reports weighted F1 before/after selection.
\begin{table}[h]
\centering
\caption{ANOVA select-$k$ summary from run metadata}
\label{tab:anova_results}
\begin{tabular}{lrrrrr}
\hline
Family & Orig. feat. & Sel. feat. & Base F1$_w$ & Sel. F1$_w$ & $\Delta$F1$_w$ \\
\hline
bert & 768 & 300 & 0.5191 & 0.5100 & -0.0091 \\
mfcc\_normalized & 1150 & 300 & 0.4024 & 0.3912 & -0.0112 \\
mfcc\_raw & 1150 & 1150 & 0.3726 & 0.3726 & +0.0000 \\
prosody\_energy & 8 & 8 & 0.2236 & 0.2236 & +0.0000 \\
prosody\_pitch & 14 & 14 & 0.2439 & 0.2439 & +0.0000 \\
ssl\_hubert & 1539 & 300 & 0.3697 & 0.3443 & -0.0254 \\
ssl\_wav2vec & 1539 & 300 & 0.3681 & 0.3467 & -0.0214 \\
tfidf & 10480 & 1000 & 0.4242 & 0.3830 & -0.0412 \\
\hline
\end{tabular}
\end{table}
\subsection{feature\_selection/anova/elbow\_curves\_anova/}
\textbf{What was done.} Elbow-curve sweeps were used to estimate practical and best-performing values of $k$ for ANOVA-ranked features.
\textbf{How it was done.} For each family, notebooks evaluate multiple top-$k$ prefixes of ranked features and save summary CSV files (elbow $k$, best $k$, and scores).
\textbf{Results.} Table~\ref{tab:elbow_results} reports exported elbow summaries.
\begin{table}[h]
\centering
\caption{Elbow-curve summary (accuracy and weighted F1)}
\label{tab:elbow_results}
\begin{tabular}{llrrrr}
\hline
Family & Metric & Elbow $k$ & Best $k$ & Elbow score & Best score \\
\hline
bert & accuracy & 90 & 686 & 0.3261 & 0.3857 \\
bert & F1$_w$ & 90 & 686 & 0.3144 & 0.3771 \\
mfcc\_raw & accuracy & 220 & 848 & 0.3342 & 0.3504 \\
mfcc\_raw & F1$_w$ & 240 & 848 & 0.3136 & 0.3416 \\
prosody\_energy & accuracy & 5 & 7 & 0.2396 & 0.2492 \\
prosody\_energy & F1$_w$ & 5 & 7 & 0.2137 & 0.2230 \\
ssl\_hubert & accuracy & 70 & 280 & 0.3519 & 0.3660 \\
ssl\_hubert & F1$_w$ & 60 & 280 & 0.3246 & 0.3542 \\
ssl\_wav2vec & accuracy & 110 & 300 & 0.3322 & 0.3569 \\
ssl\_wav2vec & F1$_w$ & 110 & 300 & 0.3138 & 0.3461 \\
tfidf & accuracy & 170 & 814 & 0.2218 & 0.2615 \\
tfidf & F1$_w$ & 170 & 814 & 0.1369 & 0.2409 \\
\hline
\end{tabular}
\end{table}
\noindent Note: current elbow artifacts are split across both \texttt{feature\_selection/anova/...} and \texttt{feature\_selection/filter/...} directories.
\subsection{feature\_test/}
\textbf{What was done.} Additional notebooks test classical ML and DNN baselines on selected feature sets.
\textbf{How it was done.} Experiments include tree-based learners, linear/logistic models, SVM variants, and an MLP baseline.
\textbf{Results.} Metrics are currently stored in notebook outputs (not standardized CSV/JSON). Table~\ref{tab:feature_test_outputs} reports coarse maxima extracted from saved outputs.
\begin{table}[h]
\centering
\caption{Coarse maxima extracted from saved feature\_test notebook outputs}
\label{tab:feature_test_outputs}
\begin{tabular}{lrr}
\hline
Notebook & Max reported accuracy & Max reported F1 \\
\hline
feature\_test/classical\_feature\_system/classical\_feature\_ensemble.ipynb & n/a & 0.4304 \\
feature\_test/DNN/mfcc\_dnn\_test.ipynb & 0.4575 & 0.4000 \\
feature\_test/TreeBased/mfcc\_boosting\_test.ipynb & 0.4706 & 0.5500 \\
feature\_test/TreeBased/mfcc\_RF\_test.ipynb & 0.5698 & 0.5500 \\
feature\_test/TreeBased/text\_tfidf\_test.ipynb & 0.5234 & 0.5700 \\
\hline
\end{tabular}
\end{table}
\noindent For thesis-grade comparability, feature\_test notebooks should be extended to export structured results with fixed data splits and confidence intervals.
\subsection{feature\_selection/anova/select\_k\_best\_anova/ (combined feature set)}
\textbf{What was done.} ANOVA feature selection was extended from single families to the merged feature table, enabling a direct test of whether a compact cross-family subset could preserve performance.
\textbf{How it was done.} The notebook \texttt{k\_best\_anova\_combined.ipynb} applies the same train/test split and class-imbalance-aware evaluation strategy used in the single-family ANOVA experiments, but on the concatenated feature matrix exported by the pipeline stage.
\textbf{Results.} Table~\ref{tab:anova_combined_results} reports the current committed result for the combined table.
\begin{table}[h]
\centering
\caption{ANOVA result for the combined feature table}
\label{tab:anova_combined_results}
\begin{tabular}{lrrrrr}
\hline
Dataset & Orig. feat. & Sel. feat. & Base acc. & Sel. acc. & $\Delta$acc. \\
\hline
combined & 4859 & 300 & 0.4287 & 0.3787 & -0.0501 \\
\hline
\end{tabular}
\end{table}
\noindent The combined-table result suggests that aggressive filter selection can remove complementary information that is useful when heterogeneous feature families are fused.
\subsection{feature\_selection/pca/select\_pca/}
\textbf{What was done.} A PCA-based dimensionality-reduction workflow was added for multiple feature families, including BERT, MFCC, prosody, SSL embeddings, and TF-IDF.
\textbf{How it was done.} Family-specific PCA notebooks compare baseline models against PCA-compressed representations using a fixed train/test split. A runner notebook (\texttt{run\_all\_pca.ipynb}) executes the full batch, and each run exports transformed train/test matrices, explained-variance summaries, PCA component definitions, and run metadata.
\textbf{Results.} Table~\ref{tab:pca_results} summarizes the current committed PCA runs.
\begin{table}[h]
\centering
\caption{PCA summary from run metadata}
\label{tab:pca_results}
\begin{tabular}{lrrrrrr}
\hline
Family & Orig. feat. & PCA feat. & Var. kept & Base F1$_w$ & PCA F1$_w$ & $\Delta$F1$_w$ \\
\hline
bert & 768 & 300 & 0.9436 & 0.4145 & 0.3836 & -0.0309 \\
mfcc\_raw & 1150 & 300 & 0.9621 & 0.3726 & 0.3362 & -0.0364 \\
prosody\_energy & 8 & 8 & 1.0000 & 0.2526 & 0.2773 & +0.0246 \\
prosody\_pitch & 14 & 14 & 1.0000 & 0.2581 & 0.2529 & -0.0052 \\
ssl\_hubert & 1539 & 300 & 0.9580 & 0.3697 & 0.3817 & +0.0120 \\
ssl\_wav2vec & 1539 & 300 & 0.9427 & 0.3681 & 0.3927 & +0.0246 \\
tfidf & 10480 & 1000 & 0.5552 & 0.4044 & 0.3733 & -0.0310 \\
\hline
\end{tabular}
\end{table}
\noindent The strongest PCA gains were observed for SSL features, especially wav2vec2, where the compressed representation improved weighted F1 from 0.3681 to 0.3927. In contrast, PCA reduced performance for BERT, MFCC, and TF-IDF, indicating that variance preservation alone did not guarantee preservation of emotion-discriminative information.
\subsection{feature\_test/classical\_feature\_system/}
\textbf{What was done.} The classical-feature experiments were extended with a systematic brute-force search over feature-family combinations in addition to the earlier all-family ensemble notebook.
\textbf{How it was done.} The notebook \texttt{feature\_family\_combinations\_bruteforce.ipynb} loads all available families through shared pipeline utilities, creates a fixed stratified 80/20 train/test split, and evaluates every combination of size 2, 3, and 4 using balanced logistic regression with imputation and scaling.
\textbf{Results.} With 10 available families, the search space contains 45 two-family combinations, 120 three-family combinations, and 210 four-family combinations, for a total of 375 evaluated combinations. The notebook is configured to export a run-level CSV and summary JSON under \texttt{feature\_test/classical\_feature\_system/artifacts/family\_combo\_bruteforce/}, but no committed result artifacts are currently present in the repository. For the thesis, this experiment should therefore be described as an implemented evaluation workflow whose result tables still need to be regenerated and frozen.
\subsection{feature\_augmentation/gan\_workflow/}
\textbf{What was done.} A dedicated workflow scaffold was added for GAN-based augmentation of tabular feature sets.
\textbf{How it was done.} The current implementation consists of a planning-oriented README, a decision log, a configuration template, and a helper script to estimate class-balancing targets before model training. The intended workflow is to train a class-conditional GAN on training features only, generate synthetic minority-class samples, export augmented training tables, and then rerun the downstream classifiers without modifying the held-out test split.
\textbf{Results.} No GAN augmentation results have been committed yet. This part of the repository should therefore be written in the thesis as planned future work or as an implementation scaffold, not as a completed experiment.
\subsection{experiments/}
\textbf{What was done.} Repository-level experiment directories (\texttt{alpha}, \texttt{beta}, \texttt{gamma}) were initialized to support structured job tracking.
\textbf{How it was done.} Each directory currently contains an \texttt{index.json} descriptor and a \texttt{jobs.json} file with an empty \texttt{TRAIN} queue and no cached jobs.
\textbf{Results.} At the current stage, these directories provide experiment-management scaffolding rather than empirical findings. They are useful to mention only briefly in a reproducibility or infrastructure paragraph.
